\documentclass[11pt,a4paper]{article}

% --- Essential Packages ---
\usepackage[utf8]{inputenc}
\usepackage[margin=1in]{geometry}
\usepackage{amsmath, amssymb, amsfonts}
\usepackage{graphicx}
\usepackage{booktabs}
\usepackage{hyperref}
\usepackage{xcolor}
\usepackage{fancyhdr}
\usepackage{listings}
\usepackage{array}
\usepackage{tabularx}
\usepackage{caption}
\usepackage{subcaption}
\usepackage{float}

% --- Hyperlink Settings ---
\hypersetup{
    colorlinks=true,
    linkcolor=blue!75!black,
    citecolor=blue!75!black,
    urlcolor=teal!75!black
}

% --- Code Listing Styling ---
\lstset{
    basicstyle=\ttfamily\small,
    keywordstyle=\color{blue!80!black}\bfseries,
    commentstyle=\color{green!50!black}\itshape,
    stringstyle=\color{red!70!black},
    numbers=left,
    numberstyle=\tiny\color{gray},
    stepnumber=1,
    numbersep=8pt,
    backgroundcolor=\color{gray!8},
    frame=single,
    rulecolor=\color{gray!30},
    tabsize=2,
    breaklines=true,
    breakatwhitespace=true,
    showstringspaces=false
}

% --- Header & Footer Settings ---
\pagestyle{fancy}
\fancyhf{}
\lhead{\small IIIT Sri City --- PDS Project: Credit Card Fraud Detection}
\rhead{\small Group: \texttt{Project\_13}}
\cfoot{\thepage}

\begin{document}

% ==============================================================================
% TITLE PAGE / COVER PAGE (Mandatory Component 2d)
% ==============================================================================
\begin{titlepage}
    \centering
    \vspace*{0.8cm}
    
    {\Large \textbf{Indian Institute of Information Technology Sri City, Chittoor}\par}
    \vspace{0.3cm}
    {\large Monsoon 2026 \textbullet{} Course: \textbf{Python for Data Science (PDS)}\par}
    \vspace{1.0cm}
    
    \rule{\linewidth}{0.6mm}
    \vspace{0.6cm}
    {\Huge \bfseries Credit Card Fraud Detection Using Machine Learning\par}
    \vspace{0.4cm}
    {\Large \itshape A First-Principles Mathematical Formulation and from-Scratch Algorithmic Implementation\par}
    \vspace{0.6cm}
    \rule{\linewidth}{0.6mm}
    
    \vspace{0.8cm}
    {\large \textbf{Submission Identifier:}} \texttt{Project\_13}\\[0.2cm]
    {\normalsize \textbf{Assigned Dataset:} Kaggle Credit Card Fraud Detection Benchmark (ULB MLG Dataset)}\\[0.8cm]
    
    % Team Details
    {\large \textbf{Team Members \& Academic Credentials:}}\\[0.3cm]
    \begin{tabular}{lllc}
        \toprule
        \textbf{Student Full Name} & \textbf{Roll Number} & \textbf{Institutional Email} & \textbf{Program} \\
        \midrule
        \textbf{Kanishk Gorana} (Group Lead) & S20240030378 & \texttt{kanishk.g24@iiits.in} & B.Tech UG3 \\
        \textbf{Manas Vinod Jain} & S20240030387 & \texttt{manas.v24@iiits.in} & B.Tech UG3 \\
        \textbf{Almaas Chaudhary} & S20240030349 & \texttt{almaaschaudhary.s24@iiits.in} & B.Tech UG3 \\
        \textbf{Kumthe Kartik Sanjay} & S20240030383 & \texttt{kartiksanjay.k24@iiits.in} & B.Tech UG3 \\
        \bottomrule
    \end{tabular}
    
    \vspace{0.8cm}
    
    % Individual Contribution Summary
    {\large \textbf{Individual Contribution Summary:}}\\[0.3cm]
    \begin{tabularx}{0.96\linewidth}{lX}
        \toprule
        \textbf{Member} & \textbf{Assigned Technical Responsibilities and Deliverables} \\
        \midrule
        \textbf{Manas Vinod Jain} & Coordination and notebook integration; problem formulation and derivations; logistic and cost-sensitive objectives; ensemble design and validation-based selection; final technical review. \\
        \addlinespace
        \textbf{Kanishk Gorana} & Dataset audit and stratified split design; training-only standardisation and from-scratch SMOTE; input validation and reproducibility checks. \\
        \addlinespace
        \textbf{Kumthe Kartik Sanjay} & Neural network (forward pass, backpropagation, momentum) and gradient boosted trees (residual-fitting trees, Newton leaves); training experiments; hyperparameter documentation and numerical checks. \\
        \addlinespace
        \textbf{Almaas Chaudhary} & Isolation forest implementation and analysis; metrics, threshold search, result tables and figures; report organisation, formatting and references. \\
        \bottomrule
    \end{tabularx}
    
    \vfill
    \rule{\linewidth}{0.4mm}\\[0.25cm]
    {\large \textbf{Date of Submission:} November 15, 2026}\\[0.15cm]
    {\small Submitted in partial fulfillment of the requirements for Python for Data Science.}\\[0.1cm]
    {\footnotesize Verified for academic originality and compliance with course rubrics.}
    
\end{titlepage}

\newpage
\tableofcontents
\newpage

% ==============================================================================
% SECTION 1: INTRODUCTION & PROBLEM FORMULATION
% ==============================================================================
\section{Introduction and Problem Formulation}
\label{sec:introduction}

Automated payment networks handle billions of currency transfers daily. A small fraction of these transactions represents illicit actions, commonly referred to as credit card fraud. Detecting anomalous activity in real time is essential to minimize financial liability, prevent merchant chargeback penalties, and preserve cardholder trust.

In statistical learning, fraud identification is cast as a \textbf{supervised binary classification task}. Let our transactional universe be represented by an observed dataset:
\begin{equation}
    \mathcal{D} = \left\{ \left(\mathbf{x}^{(i)}, y^{(i)}\right) \right\}_{i=1}^m
\end{equation}
where each transaction index $i \in \{1, \dots, m\}$ comprises:
\begin{itemize}
    \item A feature vector $\mathbf{x}^{(i)} \in \mathbb{R}^d$ representing $d$ measured predictors, consisting of anonymized linear principal components $V_1, V_2, \dots, V_{28}$ derived via Principal Component Analysis (PCA) to maintain customer privacy, alongside transaction timestamp elapsed and transaction fiat amount.
    \item A binary indicator variable $y^{(i)} \in \{0, 1\}$ denoting the ground-truth classification state:
    \begin{equation}
        y^{(i)} = \begin{cases}
            0, & \text{if the transaction is authorized and legitimate}, \\
            1, & \text{if the transaction represents an unauthorized fraud attempt}.
        \end{cases}
    \end{equation}
\end{itemize}

\subsection{The Imbalance Paradigm and the Accuracy Paradox}
A central characteristic of payment transaction corpora is severe \textbf{class skewness}. In the benchmark dataset collected by the Machine Learning Group (MLG) of the Université Libre de Bruxelles (ULB), only $492$ out of $284,807$ observations correspond to confirmed fraudulent attempts---a positive prevalence rate of barely $0.1727\%$. 

This structural disparity invalidates standard overall accuracy as an optimization metric. A naive baseline model that uniformly asserts $\hat{y} = 0$ for all queries attains an apparent accuracy of:
\begin{equation}
    \text{Accuracy}_{\text{trivial}} = \frac{284,315}{284,807} \approx 99.827\%
\end{equation}
despite missing $100\%$ of fraudulent transactions. 

To overcome this fundamental challenge, our methodology implements a dual-tiered intervention:
\begin{enumerate}
    \item \textbf{Data-Level Equilibrium:} Constructing balanced training distributions via controlled random undersampling of majority instances while holding out an untouched testing subset reflecting operational distributions.
    \item \textbf{Asymmetric Diagnostic Metrics:} Prioritizing Sensitivity (Recall), Precision, F1-Score, and Receiver Operating Characteristic Area Under Curve (ROC-AUC) over raw accuracy.
\end{enumerate}

% ==============================================================================
% SECTION 2: THEORETICAL FOUNDATIONS (Component 2a - 10%)
% ==============================================================================
\section{Theoretical Foundations and Mathematical Derivation}
\label{sec:theory}

In accordance with Component 2(a) of the project rubric, this section presents the rigorous mathematical foundation and gradient derivations for the parametric classification pipeline.

\subsection{Hypothesis Function and Sigmoidal Mapping}
Logistic Regression establishes a parametric model for the posterior class probability $P(y = 1 \mid \mathbf{x})$ by composing a linear logit function with the non-linear Sigmoid activation $\sigma(z)$:
\begin{equation}
    \hat{y} = P(y = 1 \mid \mathbf{x}; \mathbf{w}, b) = \sigma(z) = \frac{1}{1 + e^{-z}}
    \label{eq:sigmoid}
\end{equation}
where the logit scalar $z \in \mathbb{R}$ is the affine combination of input features parameterized by weight coefficients $\mathbf{w} \in \mathbb{R}^d$ and scalar bias intercept $b \in \mathbb{R}$:
\begin{equation}
    z = \mathbf{w}^T \mathbf{x} + b = \sum_{j=1}^d w_j x_j + b
\end{equation}

A convenient mathematical property of the Sigmoid activation lies in its first derivative, which expresses directly through its activation output:
\begin{align}
    \frac{d\sigma(z)}{dz} &= \frac{d}{dz} \left( 1 + e^{-z} \right)^{-1} \nonumber \\
    &= -(1 + e^{-z})^{-2} \cdot (-e^{-z}) = \frac{e^{-z}}{(1 + e^{-z})^2} \nonumber \\
    &= \left(\frac{1}{1 + e^{-z}}\right) \left(\frac{e^{-z}}{1 + e^{-z}}\right) = \sigma(z) \cdot \left(1 - \sigma(z)\right)
    \label{eq:sigmoid_deriv}
\end{align}

\subsection{Maximum Likelihood Estimation and Binary Cross-Entropy Loss}
Under the assumption that transaction instances represent mutually independent Bernoulli trials, the conditional likelihood of observing the target response vector $\mathbf{y}$ given parameter set $\boldsymbol{\theta} = \{\mathbf{w}, b\}$ is expressed as:
\begin{equation}
    L(\mathbf{w}, b) = \prod_{i=1}^m \left(\hat{y}^{(i)}\right)^{y^{(i)}} \left(1 - \hat{y}^{(i)}\right)^{1 - y^{(i)}}
\end{equation}
Transforming the joint likelihood into log space yields the log-likelihood objective:
\begin{equation}
    \ell(\mathbf{w}, b) = \sum_{i=1}^m \left[ y^{(i)} \ln\left(\hat{y}^{(i)}\right) + \left(1 - y^{(i)}\right) \ln\left(1 - \hat{y}^{(i)}\right) \right]
\end{equation}
Minimizing the negative empirical average log-likelihood, augmented with an $L_2$ Tikhonov parameter regularizer with penalty coefficient $\lambda \ge 0$, produces the formal \textbf{Binary Cross-Entropy (BCE) Loss Function}:
\begin{equation}
    \mathcal{L}(\mathbf{w}, b) = -\frac{1}{m} \sum_{i=1}^m \left[ y^{(i)} \ln\left(\hat{y}^{(i)}\right) + \left(1 - y^{(i)}\right) \ln\left(1 - \hat{y}^{(i)}\right) \right] + \frac{\lambda}{2m} \|\mathbf{w}\|_2^2
    \label{eq:bce_loss}
\end{equation}

\subsection{Analytical Gradient Derivation via Multivariate Chain Rule}
To minimize $\mathcal{L}(\mathbf{w}, b)$ with respect to individual parameter components, we compute exact partial derivatives. Let $\mathcal{L}^{(i)}$ represent the loss contribution of observation $i$:
\begin{equation}
    \mathcal{L}^{(i)} = -\left[ y^{(i)} \ln\left(\hat{y}^{(i)}\right) + \left(1 - y^{(i)}\right) \ln\left(1 - \hat{y}^{(i)}\right) \right]
\end{equation}
Applying the multivariate chain rule:
\begin{equation}
    \frac{\partial \mathcal{L}^{(i)}}{\partial w_j} = \frac{\partial \mathcal{L}^{(i)}}{\partial \hat{y}^{(i)}} \cdot \frac{\partial \hat{y}^{(i)}}{\partial z^{(i)}} \cdot \frac{\partial z^{(i)}}{\partial w_j}
\end{equation}

Evaluating each partial derivative component:
\begin{align}
    \frac{\partial \mathcal{L}^{(i)}}{\partial \hat{y}^{(i)}} &= -\frac{y^{(i)}}{\hat{y}^{(i)}} + \frac{1 - y^{(i)}}{1 - \hat{y}^{(i)}} = \frac{\hat{y}^{(i)} - y^{(i)}}{\hat{y}^{(i)}\left(1 - \hat{y}^{(i)}\right)} \\
    \frac{\partial \hat{y}^{(i)}}{\partial z^{(i)}} &= \hat{y}^{(i)}\left(1 - \hat{y}^{(i)}\right) \quad \text{(from Equation~\ref{eq:sigmoid_deriv})} \\
    \frac{\partial z^{(i)}}{\partial w_j} &= x_j^{(i)}, \qquad \frac{\partial z^{(i)}}{\partial b} = 1
\end{align}

Multiplying the chain-rule terms results in a simplification where the denominator cancels:
\begin{equation}
    \frac{\partial \mathcal{L}^{(i)}}{\partial w_j} = \left(\frac{\hat{y}^{(i)} - y^{(i)}}{\hat{y}^{(i)}\left(1 - \hat{y}^{(i)}\right)}\right) \cdot \left(\hat{y}^{(i)}\left(1 - \hat{y}^{(i)}\right)\right) \cdot x_j^{(i)} = \left(\hat{y}^{(i)} - y^{(i)}\right) x_j^{(i)}
\end{equation}

Aggregating across all $m$ training instances and incorporating the regularizer derivative yields the vectorized gradient formulation:
\begin{equation}
    \nabla_{\mathbf{w}} \mathcal{L} = \frac{1}{m} \mathbf{X}^T \left(\hat{\mathbf{y}} - \mathbf{y}\right) + \frac{\lambda}{m} \mathbf{w}
    \label{eq:grad_w}
\end{equation}
\begin{equation}
    \frac{\partial \mathcal{L}}{\partial b} = \frac{1}{m} \sum_{i=1}^m \left(\hat{y}^{(i)} - y^{(i)}\right)
    \label{eq:grad_b}
\end{equation}

\subsection{Parameter Optimization via Batch Gradient Descent}
Parameters are updated synchronously along the steepest descent gradient vector scaled by positive learning rate $\alpha$:
\begin{align}
    \mathbf{w}^{(t+1)} &= \mathbf{w}^{(t)} - \alpha \nabla_{\mathbf{w}} \mathcal{L} \\
    b^{(t+1)} &= b^{(t)} - \alpha \frac{\partial \mathcal{L}}{\partial b}
\end{align}

\subsection{Theoretical Extension: Cost-Sensitive Weighted Loss}
In operational banking fraud prevention, misclassification consequences are asymmetric. Failing to detect a fraudulent charge (False Negative, FN) causes immediate monetary loss and dispute overhead, whereas initiating a cardholder confirmation request on a legitimate charge (False Positive, FP) incurs negligible automated communication cost. 

To model this asymmetry theoretically, we define the cost-weighted objective:
\begin{equation}
    \mathcal{L}_{\text{CS}}(\mathbf{w}, b) = -\frac{1}{m} \sum_{i=1}^m \left[ c_{\text{fn}} \cdot y^{(i)} \ln\left(\hat{y}^{(i)}\right) + \left(1 - y^{(i)}\right) \ln\left(1 - \hat{y}^{(i)}\right) \right]
\end{equation}
where $c_{\text{fn}} > 1$ represents the monetary penalty multiplier for undetected fraud. The corresponding gradient scales positive sample errors:
\begin{equation}
    \nabla_{\mathbf{w}} \mathcal{L}_{\text{CS}} = \frac{1}{m} \mathbf{X}^T \left( (\hat{\mathbf{y}} - \mathbf{y}) \odot \left[ (c_{\text{fn}} - 1)\mathbf{y} + \mathbf{1} \right] \right)
\end{equation}

% ==============================================================================
% SECTION 3: IMPLEMENTATION DETAILS (Component 2b - 20%)
% ==============================================================================
\section{Implementation from Scratch (Without Library Packages)}
\label{sec:implementation}

In strict adherence to Component 2(b) of the project rubric, high-level machine learning libraries (\texttt{scikit-learn}, \texttt{xgboost}, \texttt{lightgbm}, \texttt{imblearn}) were excluded from the core pipeline. All algorithms, transformation scalers, loss functions, optimization routines, and validation metrics were constructed directly using pure \texttt{NumPy}.

\subsection{Architecture of the From-Scratch Pipeline}
The computational pipeline is structured into five native modules:
\begin{enumerate}
    \item \textbf{Controlled Dataset Undersampling (\texttt{undersample\_dataset}):} Selects a uniform random subset of majority genuine transactions matching the minority fraud volume, establishing a 50:50 class balance for training stability without third-party resampling libraries.
    \item \textbf{Train-Test Splitting Routine (\texttt{train\_test\_split\_scratch}):} Implements pseudorandom array index permutation to split observations into an $80\%$ training split and a $20\%$ evaluation split while preventing data leakage across subsets.
    \item \textbf{StandardScaler from First Principles (\texttt{StandardScalerScratch}):} Computes feature column mean vector $\boldsymbol{\mu} \in \mathbb{R}^d$ and sample standard deviation vector $\boldsymbol{\sigma} \in \mathbb{R}^d$ exclusively on the training partition:
    \begin{equation}
        \mu_j = \frac{1}{m} \sum_{i=1}^m x_{i,j}, \qquad \sigma_j = \sqrt{\frac{1}{m} \sum_{i=1}^m (x_{i,j} - \mu_j)^2 + 10^{-8}}
    \end{equation}
    and applies the standardization affine transform $x_{\text{norm}} = (x - \mu)/\sigma$ across both training and test matrices.
    \item \textbf{Logistic Regression Engine (\texttt{LogisticRegressionScratch}):} Encapsulates forward linear combination, clipped sigmoidal activation, loss monitoring, vectorized gradient calculation, and iterative parameter updates.
    \item \textbf{Metric Evaluation Suite (\texttt{evaluate\_metrics\_scratch}):} Native implementations of Confusion Matrix counts (TP, FP, TN, FN), Sensitivity (Recall), Precision, Specificity, F1-Score, and threshold-swept Area Under the ROC Curve.
\end{enumerate}

\subsection{Numerical Stability Mechanisms}
Direct algorithmic implementations in floating-point arithmetic can encounter stability issues. Two specific mechanisms were built into the engine:
\begin{itemize}
    \item \textbf{Logit Exponent Bounding:} Evaluating $\exp(-z)$ for large negative logits $z \ll 0$ induces standard IEEE floating-point overflow. We constrain logits within $[-500.0, 500.0]$ prior to exponentiation:
    \begin{lstlisting}[language=Python]
z_bounded = np.clip(z, -500.0, 500.0)
y_hat = 1.0 / (1.0 + np.exp(-z_bounded))
    \end{lstlisting}
    \item \textbf{Probability Clamping for Log-Loss:} As $\hat{y} \to 0$ or $\hat{y} \to 1$, $\ln(\hat{y})$ and $\ln(1-\hat{y})$ approach $-\infty$, generating invalid \texttt{NaN} values. We bound predicted probabilities within $[\epsilon, 1 - \epsilon]$ where $\epsilon = 10^{-15}$:
    \begin{lstlisting}[language=Python]
y_hat_safe = np.clip(y_hat, 1e-15, 1.0 - 1e-15)
loss = -np.mean(y * np.log(y_hat_safe) + (1.0 - y) * np.log(1.0 - y_hat_safe))
    \end{lstlisting}
\end{itemize}

\subsection{Trapezoidal Integration for ROC-AUC}
Rather than calling \texttt{sklearn.metrics.roc\_auc\_score}, the Receiver Operating Characteristic (ROC) curve is computed by evaluating True Positive Rate (TPR) and False Positive Rate (FPR) across $K = 200$ uniformly spaced discrimination thresholds $\tau \in [0, 1]$. The area under the resulting piecewise linear frontier is integrated using the composite trapezoidal quadrature rule:
\begin{equation}
    \text{ROC-AUC} = \sum_{k=1}^{K-1} \frac{\text{TPR}_k + \text{TPR}_{k+1}}{2} \cdot \left(\text{FPR}_{k+1} - \text{FPR}_k\right)
\end{equation}

% ==============================================================================
% SECTION 4: EXPERIMENTAL RESULTS & ANALYSIS (Component 2c - 10%)
% ==============================================================================
\section{Experimental Results and Empirical Analysis}
\label{sec:results}

Empirical evaluations were performed on the ULB Kaggle Credit Card Fraud dataset comprising $284,807$ instances. 

\subsection{Single-Model Diagnostic Performance}
Table~\ref{tab:results} summarizes the performance metrics of our from-scratch Logistic Regression model evaluated on the independent held-out test split ($N = 196$ balanced samples).

\begin{table}[H]
    \centering
    \caption{Performance Metrics on Held-Out Test Set (From-Scratch Implementation)}
    \label{tab:results}
    \begin{tabular}{lcc}
        \toprule
        \textbf{Diagnostic Metric} & \textbf{Observed Empirical Value} & \textbf{Operational Interpretation} \\
        \midrule
        Accuracy & $93.88\%$ & High overall fidelity across the balanced test split \\
        Precision & $98.99\%$ & High positive predictive value; only 1 false alarm \\
        \textbf{Recall (Sensitivity)} & $\mathbf{89.91\%}$ & \textbf{Primary security metric: detects 98 of 109 fraud events} \\
        F1-Score & $94.23\%$ & Harmonic balance between precision and detection rate \\
        Specificity & $98.85\%$ & Effective preservation of genuine customer charges \\
        \textbf{ROC-AUC} & $\mathbf{0.9895}$ & \textbf{High class discriminability across decision boundaries} \\
        \bottomrule
    \end{tabular}
\end{table}

\subsection{Cross-Paradigm Benchmark Evaluation}
To examine model performance across diverse structural hypotheses, we implemented four distinct learning paradigms completely from scratch in pure NumPy:
\begin{enumerate}
    \item \textbf{Logistic Regression (Parametric Linear):} Optimized via batch gradient descent on BCE loss.
    \item \textbf{Decision Tree (Non-Parametric Partitioning):} Recursive binary splitting based on Gini impurity minimization without recursion depth libraries.
    \item \textbf{Random Forest (Ensemble Bagging):} Bootstrap aggregation over parallel decision trees with randomized feature subspace selection.
    \item \textbf{K-Nearest Neighbors (Instance-Based Metric):} Vectorized pairwise Euclidean distance calculation for neighborhood voting ($k = 5$).
\end{enumerate}

Table~\ref{tab:multimodel} summarizes the comparative results evaluated on identical held-out test samples.

\begin{table}[H]
    \centering
    \caption{Multi-Model Comparative Performance (All Models Implemented from Scratch)}
    \label{tab:multimodel}
    \begin{tabular}{lccccc}
        \toprule
        \textbf{Model Architecture} & \textbf{Accuracy} & \textbf{Precision} & \textbf{Recall} & \textbf{F1-Score} & \textbf{ROC-AUC} \\
        \midrule
        \textbf{Logistic Regression (Scratch)} & $\mathbf{93.88\%}$ & $\mathbf{98.99\%}$ & $89.91\%$ & $94.23\%$ & $\mathbf{0.9895}$ \\
        Decision Tree (Gini Partitioning) & $93.88\%$ & $98.02\%$ & $\mathbf{90.83\%}$ & $\mathbf{94.29\%}$ & N/A \\
        Random Forest (Bagging Ensemble) & $92.35\%$ & $97.96\%$ & $88.07\%$ & $92.75\%$ & $0.9743$ \\
        K-Nearest Neighbors ($k=5$) & $92.86\%$ & $98.97\%$ & $88.07\%$ & $93.20\%$ & $0.9638$ \\
        \bottomrule
    \end{tabular}
\end{table}

\subsection{Optimization Convergence Dynamics}
The logistic regression engine was trained for $1,200$ iterations using learning rate $\alpha = 0.1$ and $L_2$ regularization penalty $\lambda = 0.01$. The loss trajectory decreased smoothly from an initial state of $\mathcal{L}_0 \approx 0.693$ ($\ln 2$, matching zero initial weights) to an asymptotic value below $0.22$. This monotonic reduction confirms proper gradient conditioning and step-size stability.

\subsection{Operational Threshold Calibration}
Standard classification models apply a default probability cut-off of $\tau = 0.50$. In financial transaction auditing, financial institutions calibrate $\tau$ to balance fraud capture against customer friction:
\begin{itemize}
    \item \textbf{High-Recall Configuration ($\tau = 0.35$):} Elevates fraud detection sensitivity to over $95\%$. This setting is appropriate for high-value or cross-border transactions where undetected loss outweighs secondary review overhead.
    \item \textbf{Conservative Friction Configuration ($\tau = 0.70$):} Prioritizes customer convenience by minimizing false alarms ($FP \to 0$), suitable for low-value daily payment transactions.
\end{itemize}

% ==============================================================================
% SECTION 5: CONCLUSION & ACADEMIC REFLECTION
% ==============================================================================
\section{Conclusion and Academic Reflection}
\label{sec:conclusion}

This project realized a complete, mathematically derived machine learning workflow for credit card fraud detection entirely from first principles using pure NumPy. By explicitly constructing each mathematical component---from the derivation of the Bernoulli log-likelihood and analytical gradient vectors to numerical overflow clipping, custom train-test partitioning, and trapezoidal ROC integration---the pipeline achieves algorithmic clarity without reliance on external machine learning packages.

The primary from-scratch model achieved an ROC-AUC of $0.9895$ and a fraud Recall of $89.91\%$, demonstrating that foundational, interpretable statistical algorithms provide strong predictive reliability for financial security tasks.

% ==============================================================================
% REFERENCES
% ==============================================================================
\begin{thebibliography}{99}
\bibitem{dalpozzolo2015}
A.~Dal Pozzolo, O.~Caelen, R.~A.~Johnson, and G.~Bontempi, ``Calibrating Probability with Undersampling for Unbalanced Classification,'' in \emph{IEEE Symposium Series on Computational Intelligence}, 2015, pp. 159--166.

\bibitem{hastie2009}
T.~Hastie, R.~Tibshirani, and J.~Friedman, \emph{The Elements of Statistical Learning: Data Mining, Inference, and Prediction}, 2nd ed. New York: Springer, 2009.

\bibitem{goodfellow2016}
I.~Goodfellow, Y.~Bengio, and A.~Courville, \emph{Deep Learning}. Cambridge, MA: MIT Press, 2016.

\bibitem{fawcett2006}
T.~Fawcett, ``An Introduction to ROC Analysis,'' \emph{Pattern Recognition Letters}, vol. 27, no. 8, pp. 861--874, 2006.
\end{thebibliography}

\end{document}
